\cotefile{26}
\scanpage{2}

\begin{flushright}
$\underline{\overline{\mathrm{V}}}$\qquad Plan\\
\unsure{Notations}
\end{flushright}
% archive sheet number: 1
\clearpage
\scanpage{3}

A) \uline{Descente et passage au quotient}

\S\,1\quad Descente fidèlement plate

\S\,2\quad \deleted{Passage au quoti\unsure{ent|.} \unsure{Généralités} sur les foncteurs sur} \inserted{Topologies sur $(\mathrm{Sch})_{/S}$. Propriétés générales des foncteurs sur $(\mathrm{Sch})_{/S}$.} \margin{[Intervertir \uline{1} et \uline{2} ??]}
\par\noindent\hfill\deleted{$\unsure{(\mathrm{Sch})_{/S}|(\mathrm{Ens})_{/S}}$}

\S\,3\quad Passage au quotient plat.

\S\,4\quad Descente non plate\quad [Application \unsure{à|:} \unsure{Ram.|Rau.|Rem.} \unsure{Samuel|Serreil} \unsure{généralisé} \unsure{\ldots|---}

\inserted{Appl. à platitude d'une \unsure{adh.} \unsure{schématique} et \unsure{lemme} de \deleted{\unsure{H}}\unsure{Hironaka}.}

\S\,5\quad Passage au quotient non plat. [\unsure{Passage} \unsure{aux|au} \unsure{quot.|quotients} \unsure{cat.|catég.} \unsure{gén}]

\S\,6\quad Descente \unsure{de|des} schémas relatifs étales, ----

[\unsure{Év\textsuperscript{t}} descente des \deleted{\unsure{morphismes}}\inserted{faisceaux} étales ?? \ldots]

B) \deleted{\uline{Propriétés générales des foncteurs sur} $(\mathrm{Sch})_{/S}$.}

\deleted{\unsure{\S\,7}}\quad \uline{Méthodes \inserted{de construction} non projectives}

\deleted{\S\,7\quad Propriétés générales \unsure{de} \unsure{foncteurs}, et propriétés de permanence}\quad \deleted{\inserted{\unsure{diverses}}}

\S\,7\quad \uline{Proreprésentabilité}\unsure{:|.} (et variantes).

\tikz[baseline=(p.base),remember picture]{\node[inner sep=1pt] (p) {\S\,\deleted{\unsure{8}}\unsure{9}};}\quad
\uline{Invariants différentiels} \deleted{\unsure{de|du} \unsure{div.}} \uline{\unsure{tangts|tang.} et normaux}

\tikz[baseline=(q.base)]{\node[inner sep=1pt] (q) {\S\,10};\draw (q.center) ellipse (1.25em and .9em);}\quad
\margin{\fbox{Après \S\,12 ??}}Propriétés générales \unsure{de|des} foncteurs, \uline{propriétés de permanence}
\margin{$\left\{\begin{gathered}\text{Foncteurs \uline{\unsure{excellent|excellents}} sur S}\\ \text{\unsure{morphismes|impliquent} excellent}\\ \text{??}\end{gathered}\right.$}

\tikz[baseline=(q.base)]{\node[inner sep=1pt] (q) {\S\,\deleted{\unsure{10|11}}11};\draw (q.center) ellipse (1.65em and .9em);}\quad
% slanted
\margin{\fbox{\parbox{.90\marginparwidth}{\raggedright\unsure{Commun au|Convient au} \unsure{par.} \unsure{20|éd} ? Cas Greenberg ?}}}
Représentabilités ``élémentaires'' \unsure{des|du} type $\prod_{\deleted{\unsure{X}}Y/S}X/Y$.

\inserted{cf SGAD $\underline{\overline{\mathrm{VIII}}}$ 6 et $\underline{\overline{\mathrm{XI}}}$ 6.8.}

\tikz[baseline=(r.base),remember picture]{\node[draw,rounded corners,inner sep=1pt] (r) {\S\,\unsure{8}\deleted{\unsure{1|11}}};}
\quad Construction de Modules cohérents
\figure{\begin{tikzpicture}[remember picture,overlay]\draw[->] (r.north west) to[out=125,in=190] (p.north);\end{tikzpicture}}

\S\,12\quad \deleted{\unsure{Représentabilité}}\inserted{Foncteurs} \unsure{repr.} par morphismes \inserted{\unsure{séparés|xprés}} loc. \unsure{pres.} \unsure{fi|finie}

\S\,13\quad Applications : des \unsure{th\textsuperscript{es}|th\textsuperscript{s}} \deleted{de} représentabilité relatifs. (y compris \unsure{(}pour les \unsure{esp.|es} foncteurs étales\ldots\unsure{,|)}

% archive sheet number: 2
\clearpage
\scanpage{4}
\mspage{(32}

% One long pencil cancellation descends from the right of the opening lines
% to the left of the last paragraph. Local obliterations are nested below.
\newsavebox{\coteFourTypescript}
\begin{lrbox}{\coteFourTypescript}
\begin{minipage}[t]{.96\linewidth}
\llap{\handwritten{\tikz[baseline]{\draw[thin] (0,0) .. controls (.16,.2) and (.34,.37) .. (.43,.17);}}}Considérant la restriction de $\handwritten{\varphi}$ à la sous-catégorie pleine de $\overset{\handwritten{\sim}}{\mathrm{X}_{\text{ét}}}$ formée des ouverts de ce topos, ou ce qui revient au même \handwritten{(6.1.)} des ouverts $\mathrm U$ de $\mathrm X$, on trouve une application
\[
\handwritten{\varphi_{0}}:\handwritten{\mathcal U}(\mathrm X)\mathrel{\mbox{-}\mkern-2mu\mbox{-}\mkern-3mu\handwritten{\to}}\handwritten{\bigl\{}0,1\handwritten{\bigr\}}
\]
commutant aux sup quelconques et aux inf finis. (On note que pour $\mathrm U\handwritten{\in\mathcal U}(\mathrm X)$, $\handwritten{\varphi}(\mathrm U)$ est vide ou réduit un point, et on prend $\handwritten{\varphi_{0}}(\mathrm U)=0$ ou $1$ suivant qu'on est dans l'un ou l'autre cas).

On voit facilement \deleted{\unsure{(}\illegible{1}\unsure{)}}, grâce au fait que tout fermé irréductible de $\mathrm X$ a un point génerique et un seul, que $\handwritten{\varphi_{0}}$ est defini à l'aide d'un \deleted{\unsure{x|e}} unique $\mathrm x\handwritten{\in}\mathrm X$, par la condition
\[
\begin{aligned}
&\handwritten{\varphi_{0}}(\mathrm U)=1\qquad\text{sss}\qquad \mathrm x\handwritten{\unsure{\mathrm E|\in}}\mathrm U\,,\\
\text{i.e.}\qquad &\handwritten{\varphi}(\mathrm U)\ne\varnothing\quad \mathrel{\handwritten{\Leftarrow}\mkern-7mu=\mkern-7mu\handwritten{\Rightarrow}}\quad \mathrm x\handwritten{\in}\mathrm U\,.
\end{aligned}
\]
Soit $\mathrm F\handwritten{\in}\operatorname{Ob}\overset{\handwritten{\sim}}{\mathrm X_{\text{ét}}}$, alors l'image de $\mathrm F$ dans le faisceau final $\mathrm X$ est un ouvert; et comme
\[
\handwritten{\varphi}(\mathrm F)\mathrel{\mbox{-}\mkern-2mu\mbox{-}\mkern-3mu\handwritten{\to}}\handwritten{\varphi}(\mathrm U)
\]
est surj\makebox[0pt][l]{\raisebox{.95ex}{\handwritten{\scriptsize\unsure{''}}}}ectif, on voit que $\handwritten{\varphi}(\mathrm F)\ne\varnothing$ sss $\handwritten{\varphi}(\mathrm U)\ne\varnothing$ i.e. sss $\mathrm x\handwritten{\in}\mathrm U$. Soit alors $\mathrm C_{\handwritten{\varphi}}$ la catégorie des couples $(\mathrm X',\handwritten{\xi})$ où $\mathrm X'\handwritten{\in}\operatorname{Ob}\mathrm X_{\text{ét}}$, et $\handwritten{\xi\in\varphi}(\mathrm X')$. On a signalé dans \handwritten{7.9.1.} que $\mathrm C_{\handwritten{\varphi}}^{\handwritten{\unsure{o|0|\circ}}}$ est pseudo-filtrante connexe. De plus grâce à \handwritten{7.9.2.} et à ce qui précéde appliqué à $\mathrm X'$ au lieu de $\mathrm X$, si $(\mathrm X',\handwritten{\xi})\handwritten{\in}\operatorname{Ob}\mathrm X_{\text{ét.}}$, il lui est\unsure{:} associé un unique point $\mathrm x'\handwritten{\in}\mathrm X'$, tel que pour un morphisme étale $\mathrm X''\mathrel{\mbox{-}\mkern-2mu\mbox{-}\mkern-2mu\mbox{-}\mkern-3mu\handwritten{\to}}\mathrm X'$, $\handwritten{\xi}$ est dans l'image de $\handwritten{\varphi}(\mathrm X'')\mathrel{\mbox{-}\mkern-2mu\mbox{-}\mkern-3mu\handwritten{\to}}\handwritten{\varphi}(\mathrm X')$ sss $\mathrm x'$ est dans celle de $\mathrm X''$. On conclut de ceci que le\deleted{\unsure{s}} préschéma\deleted{\unsure{s}} $\underset{\handwritten{\longleftarrow}}{\lim}$ des $\mathrm X'$ suivant la catégorie $\mathrm C$ des $(\mathrm X',\handwritten{\xi})$ existe et est un localisé strict $\mathrm Z$ de $\mathrm X$, correspondant donc à un point géométrique $\handwritten{\xi}$ et un foncteur \deleted{\unsure{l|b}}\unsure{f}ibre $\handwritten{{\unsure{\mathcal E|\varepsilon}}_{\xi}}$.
\end{minipage}
\end{lrbox}
\noindent\begin{tikzpicture}
\node[inner sep=0pt,outer sep=0pt] (body) {\begingroup\renewcommand{\ULthickness}{0pt}\deleted{\usebox{\coteFourTypescript}}\endgroup};
\draw[gray,thin] (body.south west) ++(.37\linewidth,.5em) -- ([xshift=-.31\linewidth,yshift=-2.6em]body.north east);
\end{tikzpicture}
% archive sheet number: 3
\clearpage
\scanpage{5}

C) \uline{\unsure{Mèthodes|Méthodes} \inserted{de construction} projectives}

\S\,14\quad Familles \unsure{limitées}.

\S\,15\quad \inserted{Foncteurs et} Schémas de Hilbert. (+ appl. aux quotients).

\S\,16\quad \inserted{Foncteurs et} Schémas de Picard.
\margin{(\unsure{Ths|Thes} d'existence projectifs \unsure{exclusivement}\deleted{\unsure{)}}; également le cas de \unsure{gpe|gpes} \unsure{structural|structuraux} affine \unsure{commutatif} plus général que $\unsure{\mathrm G_{\mathrm m}|\mathrm G_m}$\ldots)}

\S\,17\quad Variantes de Picard.

\S\,1\deleted{\unsure{7}}8\quad Applications aux courbes

\S\,1\deleted{\unsure{8}}\unsure{9|8}\quad Formalisme des normes

\S\,\deleted{\unsure{19}}\unsure{20}\quad Invariants multiprojectifs de Mumford.

D) \uline{Méthodes \inserted{de construction} mixtes}

\S\,2\deleted{0}1\quad Cas de représentabilité de $\underline{\mathrm{Pic}}$, $\underline{\mathrm{Pic}}^{0}$, $\underline{\deleted{\unsure{\mathrm P}}\mathrm{Pic}}_{X}^{\unsure{\tau|\mathrm t}}$ \ldots

(y compris le cas de Mumford utilisant \S\,\deleted{\unsure{19}}20).

\S\,21\quad Cas de représentabilité des foncteurs $\prod_{\deleted{\unsure{X}}Y/S}X/Y$ etc.

(\unsure{ths} du type Matsumura).

\noindent\rule{\linewidth}{.3pt}

NB.\quad On suppose connues \unsure{toutes} les \unsure{gén.} \unsure{fonctorielles} (topologies, descente faisceaux, quotients) ; à rappeler au besoin dans $\unsure{0|O}_{\underline{\overline{\mathrm V}}}$.

\noindent\rule{\linewidth}{.3pt}

\uline{Manquent} :

\reversemarginpar\margin{\deleted{Il \unsure{manquera}}}\margin{dans Chap. sur schémas en groupes}
$\left\{\vphantom{\begin{matrix}a\\a\end{matrix}}\right.$ les \unsure{ths} relatifs aux VA

\margin{dans Chapitre sur Picard}
\noindent$\left\{\begin{minipage}{.83\linewidth}les \unsure{ths g\textsuperscript{enx}|ths génx} \unsure{pointus|partiels} sur Picard, et l'étude qualitative générale des schémas de Picard, (y compris \unsure{critère} de \unsure{lissité} de Mumford\ldots)

(?) \deleted{\unsure{Relation} entre}\inserted{Étude du} Picard d'un schéma projectif en termes d'un pinceau de sections hyperplanes.

Les th. sur Picard genre ``\unsure{certaines|certines|sections} \unsure{équivalences|équivalentes|équivalence}''.\end{minipage}\right.$

\margin{dans Chap. sur $\unsure{\pi_1|\pi}$}
\noindent$\left\{\begin{minipage}{.83\linewidth}Les formulations de descente à la van Kampen etc, en termes de groupe fondamental\end{minipage}\right.$
\normalmarginpar
% archive sheet number: 4
\clearpage
\scanpage{6}
\mspage{(31}

% Pencil cancellation across the typescript.
\newsavebox{\coteSixTypescript}
\begin{lrbox}{\coteSixTypescript}
\begin{minipage}[t]{.96\linewidth}
(les morphismes se définissant de la façon évidente), et on note qu'on a un homomorphisme évident
\[
\underset{(\mathrm Y,\handwritten{\xi})\handwritten{\unsure{\in|\mathrm E}}\mathrm C_{\handwritten{\varphi}}^{\handwritten{\unsure{o|0|\circ}}}}{\underset{\handwritten{\longrightarrow}}{\lim}}\ \overset{\handwritten{\scriptscriptstyle\vee}}{\mathrm Y}
\mathrel{\mbox{-}\mkern-2mu\mbox{-}\mkern-2mu\mbox{-}\mkern-2mu\mbox{-}\mkern-2mu\mbox{-}\mkern-2mu\mbox{-}\mkern-3mu\handwritten{\to}}\handwritten{\varphi}
\tag*{\llap{(\(\unsure{\underline{\overline{\ast}}|\underline{\ast}|\ast}\))\hspace{.83\linewidth}}}
\]
en remarquant que
\[
\operatorname{Hom}\bigl(\underset{\mathrm C_{\handwritten{\varphi}}}{\underset{\handwritten{\rightarrow}}{\lim}}\ \overset{\handwritten{\scriptscriptstyle\vee}}{\mathrm Y},\handwritten{\varphi}\bigr)
\handwritten{\simeq}\underset{\mathrm C_{\handwritten{\varphi}}}{\underset{\handwritten{\leftarrow}}{\lim}}\operatorname{Hom}\bigl(\overset{\handwritten{\scriptscriptstyle\vee}}{\mathrm Y},\handwritten{\varphi}\bigr)
\handwritten{\simeq}\underset{\mathrm C_{\handwritten{\varphi}}}{\underset{\handwritten{\leftarrow}}{\lim}}\handwritten{\varphi}(\overset{\handwritten{\sim}}{\mathrm Y}).
\]
Or on a un élément canonique dans $\underset{\mathrm C_{\handwritten{\varphi}}}{\underset{\handwritten{\leftarrow}}{\lim}}\handwritten{\varphi}(\overset{\handwritten{\sim}}{\mathrm Y})$, en associant à tout $(\mathrm Y,\handwritten{\xi})$ l'élément $\handwritten{\xi}$ de $\handwritten{\varphi}(\mathrm Y)$. Utilisant le fait que les $\underset{\handwritten{\leftarrow}}{\lim}$ finies existent dans $\mathrm C$, et que le foncteur $\mathrm Y\handwritten{\leadsto\varphi}(\overset{\handwritten{\sim}}{\mathrm Y})$ y commute, on voit que dans $\mathrm C_{\handwritten{\varphi}}$ les $\underset{\handwritten{\leftarrow}}{\lim}$ finies existent; à fortion $\mathrm C_{\handwritten{\varphi}}^{\handwritten{o}}$ est pseudo-filtrant. Utilisant que tout faisceau $\mathrm F$ est limite inductive des faisceaux de la forme $\overset{\handwritten{\sim}}{\mathrm Y}$, et utilisant le fait que $\handwritten{\varphi}$ commute aux limites inductives, on conclut aisement que l'homomorphisme de foncteurs ($\underline{\overline{\ast}}$) est bijectif, et donne donc la représentation annoncée.

\handwritten{7.9.2.} Remarquons également que si $\underline{\mathrm E}$ est un topos, $\mathrm Y$ un objet de $\underline{\mathrm E}$, d'où un topos $\underline{\mathrm E}_{/\mathrm Y}$, alors la donnée d'un foncteur fibre $\handwritten{\varphi}$ pour $\underline{\mathrm E}_{/\mathrm Y}$ équivaut à la donnée d'un couple $(\handwritten{\varphi},\handwritten{\xi})$, où $\handwritten{\varphi}$ est un foncteur fibre pour $\underline{\mathrm E}$, et $\handwritten{\xi\mathrel{\mathrm E}\varphi}(\mathrm Y)$. Au foncteur fibre $\handwritten{\psi}:\underline{\mathrm E}_{/\mathrm Y}\mathrel{\mbox{-}\mkern-2mu\mbox{-}\mkern-2mu\mbox{-}\mkern-3mu\handwritten{\to}}(\mathrm{Ens})$ on associe le couple $(\handwritten{\varphi},\handwritten{\xi})$, où $\handwritten{\varphi}$ est le composé $\handwritten{\varphi}(\mathrm Z)=\handwritten{\psi}(\mathrm Z\mathrm{x}\mathrm Y)$, $\handwritten{\xi\in\varphi}(\mathrm Y)=\handwritten{\psi}(\mathrm Y\mathrm{x}\mathrm Y)$ est l'image de l'unique élément de $\handwritten{\psi}(\mathrm Y)$ par le morphisme diagonal de $\mathrm Y\mathrm{x}\mathrm Y$.

\handwritten{7.9.3.} Revenons maintenant au cas du site $\mathrm C=\overset{\handwritten{\sim}}{\mathrm X_{\text{ét}}}$, et soit $\handwritten{\varphi}:\overset{\handwritten{\sim}}{\mathrm X_{\text{ét}}}\mathrel{\mbox{-}\mkern-2mu\mbox{-}\mkern-2mu\mbox{-}\mkern-3mu\handwritten{\to}}(\mathrm{Ens})$ un foncteur fibre sur le topos étale de $\mathrm X$.
\end{minipage}
\end{lrbox}
\noindent\begin{tikzpicture}
\node[inner sep=0pt,outer sep=0pt] (body) {\begingroup\renewcommand{\ULthickness}{0pt}\deleted{\usebox{\coteSixTypescript}}\endgroup};
\draw[gray,thin] ([xshift=.44\linewidth,yshift=.8em]body.south west) -- ([xshift=-.43\linewidth,yshift=-1.3em]body.north east);
\end{tikzpicture}
% archive sheet number: 5
\clearpage
\scanpage{7}

\uline{\unsure{Recommandation} importante}

Partout où c'est possible, remplacer ``\unsure{fid|fid.} \unsure{plt|plt.} \unsure{qu.|qu} \unsure{cpct|cpt}'' par ``couvrant''.

\noindent\rule{\linewidth}{.3pt}

(pré)-Relation d'équivalence en pièces \unsure{détachées} \unsure{??|!?}

\diagram{\begin{tikzpicture}
\node[inner sep=5pt] (c) {\begin{tikzcd}[ampersand replacement=\&,column sep=small]
\deleted{\unsure{X|S}}\& (R_{0}^{\deleted{\unsure{o}}i}) \arrow[l,phantom,"\deleted{\unsure{\leftarrow|=}}" description] \& (R_{1}^{ij}) \arrow[l,shift left=1.6pt] \arrow[l,shift right=1.6pt]
\end{tikzcd}};
\draw (c.center) ellipse (3.4cm and .85cm);
\end{tikzpicture}}

% Paper watermark, sideways, seen in reverse through the sheet.
% margin note written sideways
\margin{VELIN EXTRA P.J}
% archive sheet number: 6
\clearpage
\scanpage{8}

\margin{\deleted{\unsure{Pbs|Plan|Pl.}}}

\deleted{Th. de finitude \unsure{pr|pour} les schémas de Hilbert}

\deleted{Th de \unsure{Matsusaka} \unsure{pr|pour} les automorphismes d'une variété polarisée}

\deleted{Th de la \deleted{\unsure{la|le}} ``courbe générique'' [pour Picard]}

\deleted{Schéma de Hilbert d'une courbe \deleted{\unsure{sur propre|propre}} sur $S$.}

\deleted{Picard des courbes sur $S$ [``diviseur $\Theta$'' etc ?]}

\deleted{Picard d'un schéma propre non projectif sur $k$}

\deleted{Picard d'un $S$-schéma projectif à fibres réductibles}

\deleted{[p.ex. cas d'une courbe sur $S$]}

\diagram{\begin{tikzcd}[ampersand replacement=\&,column sep=.35em,row sep=1em,execute before arrows={
\draw[thin,gray] ([xshift=-.2em,yshift=-1.2em]\tikzcdmatrixname-1-1.south east) -- ([xshift=.4em,yshift=.8em]\tikzcdmatrixname-1-1.north east);
\draw[thin,gray] ([xshift=-.5em,yshift=-2em]\tikzcdmatrixname-4-5.south west) -- ([xshift=.4em,yshift=1em]\tikzcdmatrixname-1-5.north);
\draw[thin,gray] ([xshift=-.5em,yshift=-1.4em]\tikzcdmatrixname-4-5.south east) -- ([xshift=-.1em,yshift=1.5em]\tikzcdmatrixname-1-9.north);
\draw[thin,gray] ([xshift=-1.3em,yshift=-1.8em]\tikzcdmatrixname-4-5.south east) -- ([xshift=.4em,yshift=1.2em]\tikzcdmatrixname-1-13.north);} ]
\deleted{\mathrm A}\arrow[rr,phantom,"{\rightsquigarrow}" description] \& {} \& \deleted{\mathrm B}\arrow[ddr]\arrow[ddrrrr] \& {} \& \deleted{\mathrm C}\arrow[ddl]\arrow[ddrr] \& {} \& \deleted{\mathrm D}\arrow[dr]\arrow[rrrr,bend left=18] \& {} \& \deleted{\mathrm E}\arrow[dl] \& {} \& \deleted{\mathrm J} \& {} \& \deleted{\mathrm K}\\
{} \& {} \& {} \& {} \& {} \& {} \& {} \& \deleted{\mathrm F}\arrow[dl]\\
{} \& {} \& {} \& \deleted{\mathrm H}\arrow[dr] \& {} \& {} \& \deleted{\mathrm G}\arrow[dll]\\
{} \& {} \& {} \& {} \& \deleted{\mathrm I}
\end{tikzcd}}

\noindent\rule{\linewidth}{.3pt}

\deleted{$\mathcal L$\qquad \unsure{Appendice|Appendix} : Pbs ouverts.}

% Paper watermark, sideways.
% margin note written sideways
\margin{VELIN EXTRA P.J}
% archive sheet number: 7
\clearpage
\scanpage{9}
\noindent\hfill\begin{tabular}{|c}Notations\\\hline\end{tabular}

\diagram{\resizebox{\linewidth}{!}{\(\displaystyle\begin{array}{@{}l|c@{}}
\underline{\underline{\mathrm{Quot}}}_{F/{\unsure{x|X}}/{\unsure{S|\mathbb S}}}(S')=\Gamma(S',\underline{\mathrm{Quot}}(F'/X'/S'))=\mathrm{Quot}(F'/X'/S') & F\\[.5em]
\underline{\underline{\mathrm{Hilb}}}_{X/S}(S')=\Gamma(S',\underline{\mathrm{Hilb}}(X'/S'))=\mathrm{Hilb}(X'/S')\qquad-\qquad=\underline{\underline{\mathrm{Quot}}}_{\underline O_X/X/S}(S') & X\\[-.3em]
\quad\unsure{\prime} & \\[-.2em]
\underline{\underline{\mathrm{Div}\deleted{\unsure{\mathrm{eff}|\mathrm{rel}}}}}_{X/S}(S')=\Gamma(S',\underline{\mathrm{Div}\deleted{\mathrm{eff}}}(X'/S'))=\mathrm{Div}\deleted{\mathrm{eff}}(X'/S') & X\\[.5em]
\underline{\underline{\mathrm{Hom}}}_S(X,Y)(S')=\Gamma(S',\underline{\mathrm{Hom}}_{S'}(X',Y'))=\mathrm{Hom}_{S'}(X',Y')=\mathrm{Hom}_S(X\times_S S',Y) & X\\[.5em]
\underline{\underline{\mathrm{Isom}}}_S(X,Y)(S')=\Gamma(S',\underline{\unsure{\mathrm{Isom}|\mathrm{Imm}}}_{S'}(X',Y')=\mathrm{Isom}_{S'}(X',Y') & X\\[.5em]
\underline{\underline{\mathrm{Aut}}}_S(X)(\deleted{\unsure{=}}S')\deleted{\unsure{S'|=}}=\Gamma(S',\underline{\mathrm{Aut}}_{S'}(X'))=\mathrm{Aut}_{S'}(X') & X\\[.5em]
\underline{\underline{\mathrm{Imm}}}_S(X,Y)(S')=\Gamma(S',\underline{\mathrm{Imm}}_{S'}(X',Y'))=\mathrm{Imm}_{S'}(X',Y') & X
\end{array}\)}} 

\diagram{\resizebox{\linewidth}{!}{\(\displaystyle\begin{array}{@{}l|c@{}}
\left(\underbrace{\displaystyle\prod_{Y/S}X/Y}_{\text{ou}}\right)(S')=\Gamma(S',\deleted{\unsure{\mathrm{Hom}|\Gamma|\prod}}\underline\Gamma_{Y'/S'}(X'/Y'))=[\Gamma_{Y'/S'}(X'/Y')=]\Gamma(X'/Y') & Y\\
\underline{\underline\Gamma}_{Y/S}(X\deleted{\unsure{/}}/Y\unsure{)}=\underline{\underline H}^{0}_{Y/S}(X\deleted{\unsure{/}}/Y)\qquad\deleted{\unsure{4|\Delta}} & \\[.6em]
\underline{\underline{\mathrm{Pr\acute epic}}}_{X/S}(S')=\Gamma(S',\overset{\nearrow\,\text{\uuline{pré}faisceau}}{\underline{\mathrm{Pr\acute epic}}}_{X'/S'}\deleted{\unsure{X/S}})=\mathrm{Pr\acute epic}(X'/S')\,
\overset{\inserted{={\unsure{H1|H^1}}(X',\underline O_{X'}^{*})}}{\deleted{={\unsure{H^1}}(X',\underline O_{X'}^{*})=\Gamma(S',\underline H^{1}_{X'/S'}(\underline O^{*}_{X'}))}} & \deleted{\unsure{X}}\\[.9em]
\underline{\underline{\mathrm{Pr\acute e}\deleted{\unsure{p}}\mathrm{pic}'}}_{X/S}(S')=\Gamma(S',\underset{\searrow\,\text{faisceau}}{\underline{\mathrm{Pr\acute epic}}'}_{X'/S'})=\mathrm{Pr\acute epic}'(X'/S')\deleted{\unsure{=}}\quad=\Gamma(S',\underbrace{\underline H^{1}_{X'/S'}(\underline O^{*}_{X'})}_{\text{faisceau}}) & \\[.7em]
\underline{\underline{\mathrm{Pic}}}_{X/S}(S')=\Gamma(S',\underline{\mathrm{Pic}}_{X'/S'})=\mathrm{Pic}(X'/S')\quad=\widetilde{\underline{\underline{\mathrm{Pr\acute epic}}}}_{X/S}(\deleted{\unsure{X'/S'|X'}}S') & X
\end{array}\)}} 

\diagram{\resizebox{\linewidth}{!}{\(\displaystyle
\unsure{?}\left\{\begin{aligned}
 &\mathrm{Hom}(\overbrace{\textstyle\prod_{X/S}\underline F}^{\substack{\text{contravariant}\\\text{en }F\ \deleted{\unsure{\text{covar.}|\text{sur}}}}},\underline M)=H^0(\underline F\otimes_S\underline M),\quad
 \underline{\mathrm{Hom}}(\overbrace{\textstyle\prod_{X/S}\underline F}^{\substack{\text{contravariant}\\\text{en }F}},\underline M)=\underline H^0_{X/S}(\underline F\otimes_S\underline M)\quad\big|\,X\\
 &\underline V(\textstyle\prod_{X/S}\underline F)=\textstyle\prod_{X/S}V(\underline F)\qquad
 \boxed{\underset{?\nearrow}{V}(\textstyle\prod_{X/S}\underline F)=\underline{\underline{\mathrm{Hom}}}_{X/S}(F,\underline O_X)}
\end{aligned}\right.
\)}} 
\margin{Notation opportune\,???}

\diagram{\resizebox{\linewidth}{!}{\(\displaystyle\left\{\begin{aligned}
 &\hspace{9em}\deleted{\underline{\mathrm{Hom}}_{X/S}(F,G)(S')}\\[-.2em]
 &\deleted{\underline{\underline{\mathrm{Hom}}}_{X/S}(F,G)(S')=\Gamma(S',\underline{\mathrm{Hom}}_{X'/S'}(F',G'))=\mathrm{Hom}_{\underline O_{X'}}(F',G')\unsure{=|\sim|?}}\\
 &\mathrm{Hom}(\underline{\underline{\mathrm{Hom}}}_{X/S}(F,G),\underline M)=\mathrm{Hom}_{\underline O_X\deleted{\unsure{/S}}}(F,G\otimes_{\underline O_S}M)\qquad
 \underline{\mathrm{Hom}}_{\underline O_S}(\underline{\underline{\mathrm{Hom}}}_{X/S}(F,G),\underline M)=\underline{\mathrm{Hom}}_{{\unsure{X|Y}}/S}(F,G\otimes_{\unsure{S}}{\unsure{M|\underline M}}\\
 &\mathrm{Hom}_{\unsure{\underline O_{S'}}}(\underline{\underline{\mathrm{Hom}}}(F,G)\otimes_S S',\underline O_{S'})=\mathrm{Hom}_{\underline O_{X'}}(F',G')
 \end{aligned}\right.\quad\begin{array}{|c}\hline G\end{array}
\)}} 
% margin note written sideways
% Paper watermark passing through the notation table, partly concealed by ink.
\margin{\unsure{P.J.A|P.J.} VELIN EXTRA}
% archive sheet number: 8
\clearpage
\scanpage{10}
\hfill\unsure{-} \unsure{Problèmes|Problémes} ouverts
% archive sheet number: 9
\clearpage
\scanpage{11}
\uuline{Problèmes-types à résoudre}

\reversemarginpar
\margin{\tikz[baseline]{\node (a) at (0,0) {\S\,7};\draw (0,0) ellipse (.38 and .59);}}
\fbox{a} \uline{Anneaux de définition} pour des \unsure{structures} sans \unsure{automorphisme|automorphismes}

\margin{\tikz[baseline]{\node at (0,0) {\S\S\,20,21};\draw (0,0) ellipse (.84 and .30);}}
b) \unsure{Hilbert}, Picard sans hypothèse projective
\normalmarginpar\margin{[\unsure{Résultats partiels}]}\reversemarginpar

\margin{\tikz[baseline]{\node at (0,0) {Chap \deleted{\unsure{V}}$\underline{\overline{\mathrm{VI}}}$};\draw (0,0) ellipse (.81 and .39);}}
\unsure{Dual} d'un schéma abélien non projectif
\normalmarginpar\margin{\fbox{OK. par Raynaud}}\reversemarginpar

\margin{\tikz[baseline]{\node at (0,0) {\S\S\,3,4};\draw (0,0) ellipse (.75 and .42);}}
\fbox{c} Passage au quotient sans hypothèse de platitude

\noindent\uuline{\hspace{3cm}}

% slanted
\margin{\tikz[baseline,remember picture]{\node[rotate=13] (mumford) at (0,0) {\unsure{Mumford}};\draw[rotate=13] (0,0) ellipse (.81 and .29);}}
\noindent\(\left[\begin{array}{@{}l@{}}
\tikz[baseline,remember picture]{\coordinate (itemd);}\text{d) Passage au quotient par }\unsure{\mathrm{Gl}\,n|\mathrm{Gl}_n|\mathrm{Gl}\,m}\text{, opérant ``sans point fixe''}\\[.5em]
\text{e) Modules de variétés plus ou moins quelconques.}\\
\qquad\tikz[baseline,remember picture]{\coordinate (abel);}\text{Cas des schémas abéliens}
\end{array}\right.\)
\tikz[remember picture,overlay]{\draw[gray,thin](mumford.east) -- (itemd); \draw[gray,thin](mumford.south east) to[out=-82,in=178] (abel);}

% slanted
\margin{\tikz[baseline]{\node[align=left,inner sep=1pt] at (0,0) {Mumford -\\\unsure{Seshadri} -\\\unsure{Narashiman|Narasimhan}};\draw (0,0) ellipse (1.01 and .67);}
pour \unsure{fibrés|fibres} vectoriels\thinspace; cas fini\thinspace: \unsure{rédiger|vérifier|vérifié}}
\noindent\(\left[\begin{array}{@{}l@{}}
\text{f) Classification de fibrés principaux }\inserted{\text{sur }X/Y}\ \text{(\unsure{par modules|par un autre})}\\
\quad\text{\unsure{à} groupe affine }G\text{ sur }Y\text{, on distingue les cas}\\
\qquad\tikz[baseline,remember picture]{\coordinate (abelg);}\,G\text{ commutatif}\\
\qquad G\text{ non néc. commutatif}
\end{array}\right.\)

\margin{\tikz[baseline,remember picture]{\node (seventeen) at (0,0) {\S\,17};\draw[rotate=14] (0,0) ellipse (.53 and .24);}}
\tikz[baseline,remember picture]{\coordinate (itemg);}g) Cas particulier\thinspace; \(G\) fini

\tikz[remember picture,overlay]{\draw[gray,thin] (seventeen.east) to[out=12,in=182] (abelg); \draw[gray,thin] (seventeen.east) to[out=-21,in=168] (itemg);}
% slanted
\margin{Pb mal posé\,?}
h) Cas particulier\thinspace: \(X=P^{1}_{Y},\quad G=\unsure{\mathrm{Gl}(n)|\mathrm{Gl}(m)}\).

\noindent\uuline{\hspace{3.7cm}}

\margin{Chap $\underline{\overline{\mathrm{VI}}}$\quad\tikz[baseline]{\node at (0,0) {oui};\draw[rotate=13] (0,0) ellipse (.53 and .35);}}
i) Les modules formels d'une variété abélienne \deleted{\unsure{$A$|$S$}} sont-ils \uline{simples} sur l'anneau de base\,?

\noindent\uuline{\hspace{3.8cm}}

\fbox{j)} ``Th. du carré'' dans \unsure{cas} non projectif \(\unsure{-\, -\, -\, -|\ldots}\)

% margin note written sideways
% Paper watermark through the centre.
\margin{EXTRA \unsure{P.J.A|P.J.} VELIN}
% archive sheet number: 10
